% !TEX root = ../main.tex
%% Corrigés — Chapitre 9 — Calcul matriciel et systèmes linéaires
%% Résultats vérifiés par verifications/chapitre_09.py.

\corrigesChapitre{9}{Calcul matriciel et systèmes linéaires}
\corrigesSouples

\begin{blocLivre}{cor-9-1}
\begin{corrige}{1}{Combinaisons linéaires}
\demarche{Les opérations se font coefficient par coefficient.}
\begin{questions}
  \item $A+B=\begin{pmatrix}1&3\\5&4\end{pmatrix}$ ;\quad
    $2A-3B=\begin{pmatrix}2&1\\0&-17\end{pmatrix}$ ;\quad
    $A-B=\begin{pmatrix}1&1\\1&-6\end{pmatrix}$.
  \item $A=E_{1,1}+2E_{1,2}+3E_{2,1}-E_{2,2}$ : les coefficients de la
    combinaison sont ceux de la matrice.
\end{questions}
\end{corrige}
\end{blocLivre}

\begin{blocLivre}{cor-9-2}
\begin{corrige}{2}{Produits matriciels}
\demarche{Ligne contre colonne — et toujours vérifier la compatibilité des tailles.}
\begin{questions}
  \item $AB=\begin{pmatrix}4&11\\-2&-2\end{pmatrix}$ et
    $BA=\begin{pmatrix}3&-1\\17&-1\end{pmatrix}$ :
    $\res{AB\neq BA}$, le produit matriciel n'est pas commutatif.
  \item $CD=\begin{pmatrix}1&7\\2&-3\end{pmatrix}$ ($2\times2$) ; $DC$ a
    aussi un sens ($3\times2$ par $2\times3$) mais donne une matrice
    $3\times3$ — même les tailles diffèrent.
  \item $AX=\begin{pmatrix}0\\7\end{pmatrix}
    =2\begin{pmatrix}1\\3\end{pmatrix}-\begin{pmatrix}2\\-1\end{pmatrix}$ :
    c'est bien $2\,C_1-C_2$, combinaison des colonnes de $A$ par les
    coefficients de $X$.
\end{questions}
\end{corrige}
\end{blocLivre}

\begin{blocLivre}{cor-9-3}
\begin{corrige}{3}{Matrices élémentaires et symbole de Kronecker}
\demarche{Le produit des $E_{i,j}$ est nul sauf raccord des indices.}
\begin{questions}
  \item $E_{1,2}E_{2,3}=\res{E_{1,3}}$ et $E_{1,2}E_{1,2}=\res{0}$.
  \item $(E_{i,j}E_{k,l})_{r,s}=\sum_{t}\delta_{r,i}\delta_{j,t}\,
    \delta_{t,k}\delta_{l,s}=\delta_{r,i}\,\delta_{j,k}\,\delta_{l,s}$ :
    c'est le coefficient $(r,s)$ de $\delta_{j,k}E_{i,l}$.
  \item $E_{1,2}$ et $E_{1,2}$ (ou $E_{1,2}$ et $E_{3,1}$\dots) : deux
    matrices non nulles de produit nul — l'anneau des matrices n'est pas
    intègre.
\end{questions}
\end{corrige}
\end{blocLivre}

\begin{blocLivre}{cor-9-4}
\begin{corrige}{4}{Transposées, symétriques, antisymétriques}
\demarche{La moyenne avec la transposée symétrise ; la demi-différence antisymétrise.}
\begin{questions}
  \item $(AB)^{\top}=\begin{pmatrix}4&-2\\11&-2\end{pmatrix}
    =B^{\top}A^{\top}$ : conforme.
  \item $A=\underbrace{\begin{pmatrix}1&5/2\\5/2&-1\end{pmatrix}}_{S}
    +\underbrace{\begin{pmatrix}0&-1/2\\1/2&0\end{pmatrix}}_{K}$,
    avec $S^{\top}=S$ et $K^{\top}=-K$.
  \item Si $A=S+K=S'+K'$ : en transposant, $A^{\top}=S-K=S'-K'$ ; par
    somme et différence, $S=S'=\frac{A+A^{\top}}{2}$ et
    $K=K'=\frac{A-A^{\top}}{2}$ : unicité.
\end{questions}
\end{corrige}
\end{blocLivre}

\begin{blocLivre}{cor-9-5}
\begin{corrige}{5}{Premières puissances}
\demarche{Conjecturer sur les petits exposants, prouver par récurrence.}
\begin{questions}
  \item $D^{n}=\operatorname{diag}\bigl(2^{n},(-3)^{n}\bigr)$.
  \item $T^{2}=\begin{pmatrix}1&2\\0&1\end{pmatrix}$,
    $T^{3}=\begin{pmatrix}1&3\\0&1\end{pmatrix}$ : on conjecture
    $\res{T^{n}=\begin{pmatrix}1&n\\0&1\end{pmatrix}}$. Hérédité :
    $T^{n+1}=T^{n}T$ a pour coefficient $(1,2)$ la valeur $n+1$.
  \item $J^{2}=2J$, puis par récurrence
    $J^{n}=2^{n-1}J$ : si $J^{n}=2^{n-1}J$, alors
    $J^{n+1}=2^{n-1}J^{2}=2^{n}J$.
\end{questions}
\end{corrige}
\end{blocLivre}

\begin{blocLivre}{cor-9-6}
\begin{corrige}{6}{Inversibles $2\times2$}
\demarche{Le produit avec la « comatrice » fait apparaître $(ad-bc)\,I_2$.}
\begin{questions}
  \item Un calcul direct donne
    $A\begin{pmatrix}d&-b\\-c&a\end{pmatrix}
    =\begin{pmatrix}d&-b\\-c&a\end{pmatrix}A=(ad-bc)\,I_2$.
    Si $ad-bc\neq0$, on divise : $A$ est inversible avec la formule
    annoncée. Si $ad-bc=0$ : $AB=0$ avec
    $B=\begin{pmatrix}d&-b\\-c&a\end{pmatrix}$ ; si $A$ était inversible,
    $B=0$, donc $A=0$ : absurde (sauf $A=0$, non inversible).
  \item $2\times3-5\times1=1\neq0$ :
    $\res{\begin{pmatrix}2&5\\1&3\end{pmatrix}^{-1}
    =\begin{pmatrix}3&-5\\-1&2\end{pmatrix}}$.
    Pour la seconde : $1\times4-2\times2=0$ : \res{\text{non inversible}}.
\end{questions}
\end{corrige}
\end{blocLivre}

\begin{blocLivre}{cor-9-7}
\begin{corrige}{7}{Un système, vu matriciellement}
\demarche{Résoudre, puis relire la solution comme des coefficients de combinaison.}
\begin{questions}
  \item $A=\begin{pmatrix}1&1&1\\1&2&3\\2&1&-1\end{pmatrix}$,
    $X=\begin{pmatrix}x\\y\\z\end{pmatrix}$,
    $B=\begin{pmatrix}6\\14\\1\end{pmatrix}$.
  \item La solution est $(1,2,3)$ (chapitre 2). Donc
    $B=1\,C_1+2\,C_2+3\,C_3$ : le second membre est combinaison linéaire
    des colonnes de $A$ — c'est exactement la compatibilité du système.
\end{questions}
\end{corrige}
\end{blocLivre}

\begin{blocLivre}{cor-9-8}
\begin{corrige}{8}{Un inverse par la résolution de $AX=Y$}
\demarche{Résoudre en gardant un second membre littéral : la solution lit l'inverse.}
\begin{questions}
  \item $x+2y=u$ et $x+3y=v$ : par différence, $y=v-u$, puis
    $x=u-2(v-u)=3u-2v$.
  \item La solution s'écrit
    $X=\begin{pmatrix}3&-2\\-1&1\end{pmatrix}Y$ pour tout $Y$ : $A$ est
    inversible et
    $\res{A^{-1}=\begin{pmatrix}3&-2\\-1&1\end{pmatrix}}$ ;
    on vérifie $AA^{-1}=I_2$.
\end{questions}
\end{corrige}
\end{blocLivre}

\begin{blocLivre}{cor-9-9}
\begin{corrige}{9}{Diviseurs de zéro et nilpotents}
\demarche{Les matrices élémentaires fournissent tous les contre-exemples voulus.}
\begin{questions}
  \item $A=\begin{pmatrix}0&1\\0&0\end{pmatrix}$ et
    $B=\begin{pmatrix}1&0\\0&0\end{pmatrix}$ : $AB=0$ tandis que
    $BA=\begin{pmatrix}0&1\\0&0\end{pmatrix}\neq0$.
  \item $N^{2}=0$ : $N$ est nilpotente. Le chapitre 8 donne alors
    $(I_2+N)^{-1}=I_2-N=\begin{pmatrix}1&-1\\0&1\end{pmatrix}$
    (la somme géométrique s'arrête à $N$).
  \item Si $N^{p}=0$ et $N$ inversible, alors
    $0=N^{p}(N^{-1})^{p}=I_n$ : absurde. Une nilpotente n'est jamais
    inversible.
\end{questions}
\end{corrige}
\end{blocLivre}

\begin{blocLivre}{cor-9-10}
\begin{corrige}{10}{Le binôme matriciel}
\demarche{$I$ commute avec tout : le binôme s'applique, et la nilpotence tronque la somme.}
\begin{questions}
  \item $N^{2}=E_{1,3}$ et $N^{3}=0$.
  \item $I_3$ et $N$ commutent, donc
    $A^{n}=\sum_{k}\binom{n}{k}N^{k}=I_3+nN+\binom{n}{2}N^{2}$ (les termes
    suivants sont nuls) :
    \[
      \res{A^{n}=\begin{pmatrix}1&n&\binom{n}{2}\\0&1&n\\0&0&1\end{pmatrix}}.
    \]
  \item $n=2$ : $\begin{pmatrix}1&2&1\\0&1&2\\0&0&1\end{pmatrix}=A^{2}$ : conforme ;
    $n=3$ : coefficient $(1,3)$ égal à $3$ : conforme.
\end{questions}
\end{corrige}
\end{blocLivre}

\begin{blocLivre}{cor-9-11}
\begin{corrige}{11}{Puissances : deux méthodes}
\demarche{Écrire $\lambda I+N$ : le binôme ne garde que deux termes.}
\begin{questions}
  \item $A=2I_2+N$ avec $N=\begin{pmatrix}0&1\\0&0\end{pmatrix}$,
    $N^{2}=0$ :
    \[
      A^{n}=2^{n}I_2+n\,2^{n-1}N
      =\res{\begin{pmatrix}2^{n}&n\,2^{n-1}\\0&2^{n}\end{pmatrix}}.
    \]
  \item \etape{Initialisation} $n=1$ : immédiate.
    \etape{Hérédité} $A^{n+1}=A^{n}A$ a pour coefficient $(1,2)$ :
    $2^{n}\cdot1+n2^{n-1}\cdot2=(n+1)2^{n}$, et pour coefficients
    diagonaux $2^{n+1}$ : c'est la formule au rang $n+1$.
\end{questions}
\end{corrige}
\end{blocLivre}

\begin{blocLivre}{cor-9-12}
\begin{corrige}{12}{Suites couplées}
\demarche{Vectoriser la récurrence : les puissances de la matrice font tout le travail.}
\begin{questions}
  \item $X_{n+1}=\begin{pmatrix}2&1\\0&2\end{pmatrix}X_n$, et une
    récurrence immédiate donne $X_n=A^{n}X_0$.
  \item Avec $X_0=\begin{pmatrix}1\\1\end{pmatrix}$ et l'exercice 11 :
    \[
      \res{u_n=2^{n}+n\,2^{n-1}\qquad v_n=2^{n}}.
    \]
    (Contrôle : $u_1=2+1=3=2u_0+v_0$.)
\end{questions}
\end{corrige}
\end{blocLivre}

\begin{blocLivre}{cor-9-13}
\begin{corrige}{13}{Symétriques : produits et pièges}
\demarche{Transposer le produit et comparer : tout est dans $(AB)^{\top}=BA$.}
\begin{questions}
  \item $(AA^{\top})^{\top}=(A^{\top})^{\top}A^{\top}=AA^{\top}$ :
    symétrique (de même $A^{\top}A$).
  \item $S_1=\begin{pmatrix}1&2\\2&0\end{pmatrix}$ et
    $S_2=\begin{pmatrix}0&1\\1&3\end{pmatrix}$ sont symétriques, mais
    $S_1S_2=\begin{pmatrix}2&7\\0&2\end{pmatrix}$ ne l'est pas.
  \item Pour $A,B$ symétriques : $(AB)^{\top}=B^{\top}A^{\top}=BA$.
    Donc $AB$ symétrique $\iff AB=(AB)^{\top}=BA$.
\end{questions}
\end{corrige}
\end{blocLivre}

\begin{blocLivre}{cor-9-14}
\begin{corrige}{14}{Matrices qui commutent}
\demarche{Écrire la matrice inconnue avec quatre coefficients et identifier.}
\begin{questions}
  \item Pour $X=\begin{pmatrix}a&b\\c&d\end{pmatrix}$ :
    $DX=\begin{pmatrix}a&b\\2c&2d\end{pmatrix}$ et
    $XD=\begin{pmatrix}a&2b\\c&2d\end{pmatrix}$ : l'égalité impose
    $b=c=0$. Le commutant de $D$ est l'ensemble des
    \res{\text{matrices diagonales}}.
  \item $NX=\begin{pmatrix}c&d\\0&0\end{pmatrix}$ et
    $XN=\begin{pmatrix}0&a\\0&c\end{pmatrix}$ : $c=0$ et $a=d$. Le
    commutant de $N$ est $\res{\ens{aI_2+bN \mid a,b\in\K}}$.
  \item Si $XM=MX$ et $YM=MY$ : $(\lambda X+\mu Y)M=M(\lambda X+\mu Y)$
    et $(XY)M=XMY=M(XY)$ : le commutant est stable par combinaison
    linéaire et par produit (c'est un sous-anneau).
\end{questions}
\end{corrige}
\end{blocLivre}

\begin{blocLivre}{cor-9-15}
\begin{corrige}{15}{Gauss-Jordan : un inverse $3\times3$}
\demarche{Les opérations qui mènent $A$ à $I_3$ mènent $I_3$ à $A^{-1}$.}
\begin{questions}
  \item Sur $(A\mid I_3)$ : $L_3\leftarrow L_3-L_1$, puis
    $L_3\leftarrow L_3-L_2$, puis $L_3\leftarrow-L_3$ échelonnent ;
    la remontée ($L_2\leftarrow L_2-2L_3$, $L_1\leftarrow L_1-L_3$,
    $L_1\leftarrow L_1-L_2$) donne
    \[
      \res{A^{-1}=\begin{pmatrix}2&0&-1\\-2&-1&2\\1&1&-1\end{pmatrix}}.
    \]
  \item $AA^{-1}=I_3$ : vérification directe.
  \item $X=A^{-1}\begin{pmatrix}2\\1\\3\end{pmatrix}
    =\res{\begin{pmatrix}1\\1\\0\end{pmatrix}}$
    (contrôle : $1+1+0=2$, $0+1+2\cdot0=1$, $1+2+0=3$).
\end{questions}
\end{corrige}
\end{blocLivre}

\begin{blocLivre}{cor-9-16}
\begin{corrige}{16}{Matrices triangulaires}
\demarche{Dans un produit triangulaire, seuls les chemins d'indices croissants survivent.}
\begin{questions}
  \item Pour $i>j$ :
    $(TT')_{i,j}=\sum_{k}t_{i,k}t'_{k,j}$, et chaque terme exige $k\geq i$
    (sinon $t_{i,k}=0$) et $k\leq j$ (sinon $t'_{k,j}=0$) : impossible si
    $i>j$. Pour $i=j$, seul $k=i$ survit :
    $(TT')_{i,i}=t_{i,i}t'_{i,i}$.
  \item Si tous les $t_{i,i}\neq0$, le système $TX=Y$ se résout de bas en
    haut, avec une solution unique pour chaque $Y$ : $T$ est inversible.
    Si un $t_{i,i}=0$, le système homogène possède une solution non
    nulle (on remonte en choisissant $x_i=1$) : $T$ n'est pas inversible.
  \item Résolution de $TX=Y$ en remontant :
    \[
      \res{T^{-1}=\begin{pmatrix}1&-\tfrac12&-\tfrac{1}{12}\\[1mm]
      0&\tfrac14&-\tfrac{5}{24}\\[1mm]0&0&\tfrac16\end{pmatrix}},
    \]
    triangulaire supérieure, de diagonale $\bigl(1,\tfrac14,\tfrac16\bigr)$
    — les inverses des coefficients diagonaux.
\end{questions}
\end{corrige}
\end{blocLivre}

\begin{blocLivre}{cor-9-17}
\begin{corrige}{17}{Un paramètre dans le système}
\demarche{Le cas inversible se résout d'un coup ; le cas dégénéré se discute.}
\begin{questions}
  \item $ad-bc=m-1$ : $A_m$ est inversible si et seulement si
    $\res{m\neq1}$. Alors, par différence des lignes :
    $y=\dfrac{1}{m-1}$ et $x=2-\dfrac{1}{m-1}=\dfrac{2m-3}{m-1}$.
  \item Pour $m=1$, le système devient $x+y=2$ et $x+y=3$ :
    \res{\text{incompatible}}. Le système $x+y=2$ (deux fois) est
    compatible, de solutions
    \[
      \res{X=\begin{pmatrix}2\\0\end{pmatrix}
      +t\begin{pmatrix}1\\-1\end{pmatrix},\quad t\in\R} :
    \]
    une solution particulière plus les solutions de l'homogène
    $x+y=0$.
\end{questions}
\end{corrige}
\end{blocLivre}

\begin{blocLivre}{cor-9-18}
\begin{corrige}{18}{Matrices de rotation}
\demarche{Les formules d'addition du chapitre 2, rangées dans une matrice.}
\begin{questions}
  \item Le produit $R(a)R(b)$ a pour première colonne
    $\begin{pmatrix}\cos a\cos b-\sin a\sin b\\
    \sin a\cos b+\cos a\sin b\end{pmatrix}
    =\begin{pmatrix}\cos(a+b)\\\sin(a+b)\end{pmatrix}$, et de même pour la
    seconde : $R(a)R(b)=R(a+b)$.
  \item $R(\theta)R(-\theta)=R(0)=I_2$ :
    $R(\theta)^{-1}=R(-\theta)$, qui est aussi $R(\theta)^{\top}$
    (parité de $\cos$, imparité de $\sin$). Par récurrence,
    $R(\theta)^{n}=R(n\theta)$.
  \item La question 1 dit exactement que $\theta\mapsto R(\theta)$ est un
    morphisme de $(\R,+)$ dans $\mathrm{GL}_2(\R)$. Son noyau est
    $\ens{\theta \mid \cos\theta=1,\ \sin\theta=0}=\res{2\pi\Z}$ — le même
    noyau que $t\mapsto\mathrm{e}^{\mathrm{i}t}$, et ce n'est pas un
    hasard (exercice 24).
\end{questions}
\end{corrige}
\end{blocLivre}

\begin{blocLivre}{cor-9-19}
\begin{corrige}{19}{Le centre de $\mathcal{M}_n(\K)$}
\demarche{Tester contre les matrices élémentaires suffit — et contraint tout.}
\begin{questions}
  \item Toute matrice est combinaison linéaire des $E_{i,j}$ : si $Z$
    commute avec chaque $E_{i,j}$, il commute avec toutes les matrices
    (bilinéarité du produit).
  \item $ZE_{i,j}$ a pour seule colonne non nulle la colonne $j$, égale à
    la colonne $i$ de $Z$ ; $E_{i,j}Z$ a pour seule ligne non nulle la
    ligne $i$, égale à la ligne $j$ de $Z$. L'égalité impose : tous les
    coefficients hors diagonale de $Z$ sont nuls, et
    $z_{i,i}=z_{j,j}$ pour tous $i,j$ : $Z=\lambda I_n$.
  \item Réciproquement, les matrices scalaires commutent avec tout :
    \res{\text{le centre de }\mathcal{M}_n(\K)\text{ est }\ens{\lambda I_n}}.
\end{questions}
\end{corrige}
\end{blocLivre}

\begin{blocLivre}{cor-9-20}
\begin{corrige}{20}{Une équation matricielle}
\demarche{Commuter d'abord, paramétrer ensuite, résoudre enfin.}
\begin{questions}
  \item $XT=XX^{2}=X^{2}X=TX$ : toute solution commute avec $T$, donc
    avec $N=T-I_2$.
  \item D'après l'exercice 14, $X=aI_2+bN$.
  \item $X^{2}=a^{2}I_2+2abN$ (car $N^{2}=0$), et $T=I_2+N$ : on identifie
    $a^{2}=1$ et $2ab=1$, d'où $a=\pm1$ et $b=\dfrac{1}{2a}$ :
    \[
      \res{X=\pm\begin{pmatrix}1&\tfrac12\\0&1\end{pmatrix}} :
    \]
    exactement deux « racines carrées » — et l'on note que $T$, pourtant
    très simple, n'en a que deux.
\end{questions}
\end{corrige}
\end{blocLivre}

\begin{blocLivre}{cor-9-21}
\begin{corrige}{21}{Fibonacci, vu par les matrices}
\demarche{La récurrence de Fibonacci est un produit matriciel déguisé.}
\begin{questions}
  \item \etape{Initialisation} $A^{1}=\begin{pmatrix}1&1\\1&0\end{pmatrix}
    =\begin{pmatrix}F_2&F_1\\F_1&F_0\end{pmatrix}$.
    \etape{Hérédité} Si la formule vaut au rang $n$ :
    \[
      A^{n+1}=A^{n}A=\begin{pmatrix}F_{n+1}+F_n&F_{n+1}\\
      F_n+F_{n-1}&F_n\end{pmatrix}
      =\begin{pmatrix}F_{n+2}&F_{n+1}\\F_{n+1}&F_n\end{pmatrix}.
    \]
  \item Développement direct : pour
    $M=\begin{pmatrix}a&b\\c&d\end{pmatrix}$ et
    $M'=\begin{pmatrix}a'&b'\\c'&d'\end{pmatrix}$, le « $ad-bc$ » de
    $MM'$ vaut $(ad-bc)(a'd'-b'c')$ (huit termes qui se regroupent).
  \item Le « $ad-bc$ » de $A$ vaut $-1$ ; celui de $A^{n}$ vaut donc
    $(-1)^{n}$, c'est-à-dire
    $F_{n+1}F_{n-1}-F_n^{2}=(-1)^{n}$ : l'identité de Cassini en une
    ligne.
\end{questions}
\end{corrige}
\end{blocLivre}

\begin{blocLivre}{cor-9-22}
\begin{corrige}{22}{Matrices de permutation}
\demarche{Permuter les colonnes de $I_3$, et composer.}
\begin{questions}
  \item $P_\tau=\begin{pmatrix}0&1&0\\1&0&0\\0&0&1\end{pmatrix}$ et
    $P_c=\begin{pmatrix}0&0&1\\1&0&0\\0&1&0\end{pmatrix}$ pour
    $c=(1\,2\,3)$.
  \item $(P_\sigma P_\tau)_{i,j}=\sum_k\delta_{i,\sigma(k)}\delta_{k,\tau(j)}
    =\delta_{i,\sigma(\tau(j))}$ : c'est $P_{\sigma\circ\tau}$. Le
    morphisme est injectif : $P_\sigma=I_3$ impose $\sigma=\mathrm{id}$
    (noyau trivial).
  \item $(P_\sigma^{\top}P_\sigma)_{i,j}=\sum_k\delta_{k,\sigma(i)}
    \delta_{k,\sigma(j)}=\delta_{\sigma(i),\sigma(j)}=\delta_{i,j}$ :
    donc $P_\sigma^{\top}P_\sigma=I_3$ et
    $\res{P_\sigma^{-1}=P_\sigma^{\top}}$ ($=P_{\sigma^{-1}}$).
\end{questions}
\end{corrige}
\end{blocLivre}

\begin{blocLivre}{cor-9-23}
\begin{corrige}{23}{Un modèle d'évolution : ville et campagne}
\demarche{Deux combinaisons bien choisies découplent le système.}
\begin{questions}
  \item $u_{n+1}=0{,}9u_n+0{,}2v_n$ et $v_{n+1}=0{,}1u_n+0{,}8v_n$ :
    c'est $X_{n+1}=AX_n$. En sommant :
    $s_{n+1}=u_{n+1}+v_{n+1}=u_n+v_n=s_n$ : la population totale reste
    $\res{s_n=60}$.
  \item $d_{n+1}=u_{n+1}-2v_{n+1}=0{,}7u_n-1{,}4v_n=0{,}7\,d_n$ :
    géométrique de raison $0{,}7$, avec $d_0=30-60=-30$.
  \item De $u_n+v_n=60$ et $u_n-2v_n=-30\times0{,}7^{n}$ :
    \[
      \res{u_n=40-10\times0{,}7^{n}\qquad v_n=20+10\times0{,}7^{n}},
    \]
    de limites $40$ et $20$ (milliers d'habitants).
  \item $A\begin{pmatrix}40\\20\end{pmatrix}
    =\begin{pmatrix}36+4\\4+16\end{pmatrix}
    =\begin{pmatrix}40\\20\end{pmatrix}$ : l'état limite est bien un
    état d'équilibre, solution de $AX=X$ — le système « oublie » sa
    répartition initiale et converge vers l'équilibre.
\end{questions}
\end{corrige}
\end{blocLivre}

\begin{blocLivre}{cor-9-24}
\begin{corrige}{24}{Les complexes, cachés dans $\mathcal{M}_2(\R)$}
\demarche{Comparer le produit matriciel au produit complexe : ils coïncident terme à terme.}
\begin{questions}
  \item $M(a,b)+M(c,d)=M(a+c,b+d)$ et
    \[
      M(a,b)\,M(c,d)=\begin{pmatrix}ac-bd&-(ad+bc)\\ad+bc&ac-bd\end{pmatrix}
      =M(ac-bd,\ ad+bc) :
    \]
    ce sont les formules de la somme et du produit de
    $(a+\mathrm{i}b)(c+\mathrm{i}d)$ !
  \item $\mathcal{C}$ contient $I_2=M(1,0)$, est stable par différence et
    par produit (question 1) : sous-anneau, \emph{commutatif} puisque le
    produit complexe l'est. L'application $a+\mathrm{i}b\mapsto M(a,b)$
    est bijective, additive, multiplicative et envoie $1$ sur $I_2$ :
    isomorphisme d'anneaux de $\Cx$ sur $\mathcal{C}$.
  \item $M(0,1)^{2}=M(-1,0)=-I_2$ : la matrice $M(0,1)$ est l'image de
    $\mathrm{i}$ — « $\mathrm{i}$ est une rotation d'un quart de tour ».
  \item Le critère $ad-bc$ donne ici $a^{2}+b^{2}$ : $M(a,b)$ est
    inversible si et seulement si $(a,b)\neq(0,0)$ — fidèle au fait que
    $\Cx$ est un corps. Enfin
    $M(\cos\theta,\sin\theta)=R(\theta)$ : la forme trigonométrique
    $r\mathrm{e}^{\mathrm{i}\theta}$ devient « homothétie de rapport $r$
    composée avec rotation d'angle $\theta$ » — la géométrie du chapitre 3
    retrouvée dans les matrices.
\end{questions}
\end{corrige}
\end{blocLivre}

% !TEX root = ../../main.tex
%% Corrigé du défi du chapitre 9 — vérifié par verifications/defis_06_10.py.

\begin{blocLivre}{cor-9-25a}
\begin{corrige}{25}{Carrés magiques et carré de Lo Shu}
\demarche{Huit équations linéaires : on les résout par substitutions
successives à partir du centre ; puis le produit par $J$ code les sommes
de lignes et de colonnes.}
\partieExo{A}{Résoudre le système des carrés magiques}
\begin{questions}
  \item Les conditions « telle somme vaut $s$ » sont linéaires : si $M$ et
    $N$ sont magiques de sommes $s$ et $t$, $\lambda M+\mu N$ est magique
    de somme $\lambda s+\mu t$. Toutes les sommes de $J$ valent $\res{3}$.
  \item Les quatre lignes passant par le centre totalisent $4s$ ; elles
    couvrent le centre quatre fois et chacune des huit autres cases une
    fois, soit la somme de tous les coefficients ($3s$, somme des trois
    lignes) plus $3m_{22}$. D'où $4s=3s+3m_{22}$ et
    $\res{m_{22}=\tfrac s3}$.
  \item Chaque ligne, colonne et diagonale de $L(c,a,b)$ contient un terme
    $c$ « central » et des termes en $\pm a$, $\pm b$ qui se compensent :
    par exemple $(c-a)+(c+a+b)+(c-b)=3c$, et de même pour les sept autres
    sommes. $L(c,a,b)$ est \res{\text{magique de somme }3c}.
  \item Avec $s=3c$ : la diagonale donne $m_{33}=2c-m_{11}=c+a$ ;
    l'antidiagonale $m_{31}=2c-m_{13}=c+b$ ; la première ligne
    $m_{12}=3c-(c-a)-(c-b)=c+a+b$ ; la première colonne
    $m_{21}=3c-(c-a)-(c+b)=c+a-b$ ; la colonne du milieu
    $m_{32}=2c-m_{12}=c-a-b$ ; la ligne du milieu
    $m_{23}=2c-m_{21}=c-a+b$. Les deux équations restantes sont alors
    automatiquement vérifiées : $\res{M=L(c,a,b)}$. Les matrices magiques
    forment une famille à \res{3} paramètres libres (un espace vectoriel
    de dimension $3$, dira le chapitre 12).
\end{questions}
\partieExo{B}{Le carré de Lo Shu}
\begin{questions}[resume]
  \item La somme des neuf coefficients vaut $45=3s$, d'où $s=15$ et
    $c=5$. Les huit autres coefficients sont $5\pm a$, $5\pm b$,
    $5\pm(a+b)$, $5\pm(a-b)$ et doivent former
    $\ens{1,2,3,4,6,7,8,9}$ : les nombres $u=|a|$, $v=|b|$,
    $|a+b|$ et $|a-b|$ — c'est-à-dire $u$, $v$, $u+v$, $|u-v|$ — forment
    $\ens{1,2,3,4}$. Le plus grand est $u+v=4$, et $u\neq v$ (sinon
    $|u-v|=0$) : $\ens{u,v}=\ens{1,3}$, $|u-v|=2$. Choix de $a=\pm1$ ou
    $\pm3$ et de $b$ correspondant : $2\times2\times2=\res{8}$ carrés.
    $L(5,1,3)=\begin{pmatrix}4&9&2\\3&5&7\\8&1&6\end{pmatrix}$ est le Lo Shu.
  \item Transposer $L(c,a,b)$ échange $c+a+b$ et $c+a-b$, $c-b$ et $c+b$,
    $c-a+b$ et $c-a-b$ : c'est $L(c,a,-b)$. Les symétries du carré
    envoient lignes sur lignes ou colonnes, diagonales sur diagonales :
    elles transforment un carré magique en carré magique. Le Lo Shu ayant
    ses quatre coins distincts, ses huit images sont distinctes : ce sont
    \res{\text{les huit carrés magiques}} à coefficients $1,\dots,9$. Il
    n'y en a « qu'un seul », à symétrie près.
\end{questions}
\end{corrige}
\end{blocLivre}

\begin{blocLivre}{cor-9-25b}
\begin{corrigeSuite}
\partieExo{C}{Calcul matriciel}
\begin{questions}[start=7]
  \item Le coefficient $(i,j)$ de $MJ$ est $\sum_k m_{ik}$, somme de la
    ligne $i$ ; celui de $JM$ est la somme de la colonne $j$. Donc $M$ est
    semi-magique de somme $s$ si et seulement si $MJ=JM=sJ$. Alors
    $(MN)J=M(tJ)=st\,J$ et $J(MN)=(sJ)N=st\,J$ : $MN$ est
    \res{\text{semi-magique de somme }st}. Si $M$ est inversible, $s\neq0$
    (sinon $MJ=0$ avec $J\neq0$, et $J=M^{-1}MJ=0$) ; de $MJ=sJ$ on tire
    $J=sM^{-1}J$, donc $M^{-1}J=\frac1sJ$, et de même $JM^{-1}=\frac1sJ$.
  \item Le calcul donne
    \[
    L^{2}=\begin{pmatrix}59&83&83\\ 83&59&83\\ 83&83&59\end{pmatrix},
    \]
    semi-magique de somme $225=15^{2}$ (question 7), mais
    \res{\text{pas magique}} : sa diagonale vaut $177$. Les produits ne
    préservent pas les diagonales.
  \item On vérifie $L\cdot\frac{1}{360}A=I_3$ ligne par colonne (par
    exemple $4\cdot23+9\cdot38+2\cdot(-37)=360$ et
    $4\cdot(-52)+9\cdot8+2\cdot68=0$). Toutes les lignes, colonnes et
    diagonales de la matrice $A$ ont pour somme $24$ (par exemple
    $23+8-7=24$ et $53+8-37=24$) : $L^{-1}$ est
    \res{\text{magique de somme }\tfrac{24}{360}=\tfrac{1}{15}} — en
    accord avec la question 7, et c'est un fait général en taille $3$.
\end{questions}
\end{corrigeSuite}
\end{blocLivre}

